Here I will give a brief review of some basic concepts in topology
that will be useful in understanding the Chern-Simons formulation of
2+1 dimensional gravity. It mainly follows the treatment in
\cite{Carlip:1998qg} However, it is not enough or rigorous for
understanding any steps in Witten's paper, but serve as an
introduction to many nouns he used in the paper.

\subsection{Homeomorphisms and diffeomorphisms}

Consider two topological manifold \(X\) and \(Y\). We can define a
\textbf{homeomorphism} between them:
\defi{
  A \textbf{homeomorphism} is a continuous bijection \(f: X \to Y\)
  with a continuous inverse \(f^{-1}: Y \to X\).
}
If such a homeomorphism exists, we say that \(X\) and \(Y\) are
\textbf{homeomorphic}, denoted as \(X \cong Y\). We view the two
manifold as "topologically equivalent", meaning they have the same
topological properties.

If the manifolds are differential manifolds, we can further define a
\textbf{diffeomorphism}:
\defi{
  A \textbf{diffeomorphism} is a smooth bijection \(f: X \to Y\) with
  a smooth inverse \(f^{-1}: Y \to X\).
}
If we focus one 3 dimensional manifolds, there is an importants theorem:
\thm{
  Moise’s Theorem

  For 2 \& 3 dimentional manifolds, being homeomorphic is equivalent
  to being diffeomorphic.
}
This theorem implies that in 2 and 3 dimensions, the topological
structure and the smooth structure are essentially the same. This is
not true in higher dimensions, where there exist manifolds that are
homeomorphic but not diffeomorphic (exotic spheres).

\subsection{Mapping Class Group}

\subsubsection{Isotopy}

We focus on the self-diffeo of one manifold $ \mathcal{M} $, $ f:
\mathcal{M} \to \mathcal{M} $. We can define an \textbf{Isotopy}
between two diffeomorphisms:
\defi{
  An \textbf{isotopy} between two diffeomorphisms \(f, g: \mathcal{M}
  \to \mathcal{M}\) is a continuous family of diffeomorphisms \(F(s):
  \mathcal{M} \to \mathcal{M}\) for \(s \in [0, 1]\) such that \(F(0)
  = f\) and \(F(1) = g\). If an Isotopy exists between \(f\) and
  \(g\), we say that they are \textbf{isotopic}.
}
Most diffeos considered in physics are isotopic to idenity map, yet
there for topologically non-trivial manifolds, there can be diffeos
that are not isotopic to identity.
\begin{itemize}
  \item eg. \textbf{Dehn Twist} on torus \(T^2\): Cut the torus along
    a non-contractible loop, twist one side by \(360^\circ\), and
    then re-glue. This operation cannot be continuously deformed to
    the identity map.
\end{itemize}

\subsubsection{Mapping Class Group}

We can define an equivalence relation using Isotopy and define its
equivalent class on the set of self-diffeos of \( \mathcal{M} \):
\defi{
  The \textbf{Mapping Class Group} (MCG) of a manifold \( \mathcal{M}
  \), denoted as \( \text{MCG}(\mathcal{M}) \), is the group of
  isotopy classes of self-diffeomorphisms of \( \mathcal{M} \). .

  We can prove that the equivalent class set has a group sturcture.
}
The mapping class group of 2 dimensional surfaces is well studies and
they are generated by:
\begin{itemize}
  \item \textbf{Dehn Twists} along non-contractible loops.
\end{itemize}

\subsection{Fundamental Group}
\subsubsection{Definition}
\defi{
  Homotopy Curve

  Consider a series of curves satisfying: $ \gamma:[0,1]\to
  M,\quad\gamma(0)=x_0,\quad\gamma(1)=x_0 $. If two curve $
  \gamma_0,\gamma_1 $ can be continuously deformed to each other
  while keeping the endpoints fixed, we say they are
  \textbf{homotopic}, denoted as $ \gamma_0\simeq\gamma_1 $.

  Mathematically, there exists a continuous map $
  F:[0,1]\times[0,1]\to M $ such that:
  \begin{align}
    \begin{aligned}&F(t,0)=\gamma_1(t)\\&F(t,1)=\gamma_2(t)\\&F(0,s)=F(1,s)=x_0.
    \end{aligned}
  \end{align}
}

\defi{
  Fundamental Group

  For a fixed base point $ x_0 $ in a topological space \( M \), the
  \textbf{fundamental group} \( \pi_1(M, x_0) \) is the set of
  equivalence classes of loops $ [\gamma] $ based at \( x_0 \) under
  the homotopy relation.
}
We can prove that the set of equivalence classes has a group
structure, with the group operation defined as the concatenation of loops.

\thm{
  For connected manifolds, the fundamental group is independent of
  the choice of base point up to isomorphism. Therefore, we often
  denote it simply as \( \pi_1(M) \).
}

\subsubsection{More on Properties}
\defi{
  Simply Connected

  A topological space \( M \) is said to be \textbf{simply connected}
  if it is path-connected and its fundamental group is trivial, i.e.,
  \( \pi_1(M) = \{e\} \), where \( e \) is the identity element.
}

\thm{
  Topological Invariance of Fundamental Group

  If two topological spaces \( M \) and \( N \) are homeomorphic,
  then their fundamental groups are isomorphic, i.e., \( \pi_1(M)
  \cong \pi_1(N) \).
}

\thm{
  Fully Characterization of 2D Surfaces by Fundamental Group

  The fundamental group fully characterizes the topology of closed
  2-dimensional surfaces. Two closed surfaces are homeomorphic if and
  only if their fundamental groups are isomorphic.
}

\subsection{Covering Space}

\subsubsection{Covering Space Definition}

We can unwrap the non-contractible loops of a manifold by considering
its covering space.

\defi{
  Covering Space

  A \textbf{covering space} of a topological space \( M \) is a
  topological space \( \tilde{M} \) together with a continuous
  surjective map \( p: \tilde{M} \to M \) such that
  \begin{itemize}
    \item For every point \( x \in M \), there exists an open
      neighborhood \( U \) of \( x \) such that the preimage \(
      p^{-1}(U) \) is homeomorphic to $ U \times \Lambda $, where \(
      \Lambda \) is a discrete set.
  \end{itemize}
}

\subsubsection{Topology of Covering Space}

The surjective map used by defining the covering space can induce a
homomorphism between the fundamental groups of the two spaces:
\thm{
  Induced Homomorphism on Fundamental Groups

  Let \( p: \tilde{M} \to M \) be a covering map. Then, for any base
  point \( \tilde{x}_0 \in \tilde{M} \) with \( p(\tilde{x}_0) = x_0
  \in M \), the map \( p \) induces a homomorphism \( p_*:
  \pi_1(\tilde{M}, \tilde{x}_0) \to \pi_1(M, x_0) \).
}
This homoporphic is injective but not necessarily surjective. The
image of \( p_* \) is a subgroup of \( \pi_1(M, x_0) \) called the
\textbf{characteristic subgroup} of the covering space. It only
consist of paths that are lifted to $ \tilde{M} $ to be loops instead
of open curves.

\thm{
  For a connected manifold, any subgroup of its fundamental group can
  be realized as the characteristic subgroup of some covering space.
}
Thus we can classify the covering spaces with the subgroups of the
fundamental group.
\defi{
  Universal Covering Space

  A \textbf{universal covering space} of a topological space \( M \)
  is a covering space \( \tilde{M} \) that is simply connected, i.e.,
  its fundamental group is trivial: \( \pi_1(\tilde{M}) = \{e\} \).
  which corresponds to the trivial subgroup of \( \pi_1(M) \).
}

\subsection{Quotient Space}

\subsubsection{Quotient Space Definition}

We can construct a new manifold by quotienting a manifold with a group action.
\defi{
  Quotient Space

  Given a group $ G $ acting on a manifold $ M $. The orbit of a
  point $ x\in M $ under the action of $ G $ is the set $ G\cdot
  x=\{g\cdot x|g\in G\} $. The \textbf{quotient space} \( M/G \) is
  the set of all orbits of points in \( M \) under the action of \( G
  \), equipped with the quotient topology.
}
We can define a natural projection from the points in \( M \) to
their corresponding orbits in \( M/G \). which gives a quotient
topology by defining open sets in \( M/G \) as those whose preimages
under the projection are open in \( M \).
\begin{itemize}
  \item In general a Quotient Space \textbf{May NOT be a Manifold}.
\end{itemize}
\thm{
  Quotient Space as Manifold

  If a group \( G \) acts properly discontinuously and freely on a
  manifold \( M \), then the quotient space \( M/G \) is also a manifold.
}
\begin{itemize}
  \item \textbf{freely:} For any point \( x \in M \), if \( g \cdot x
    = x \) for some \( g \in G \), then \( g \) must be the identity
    element of \( G \). Equivanlently the isotropy group \( G_x = \{g
    \in G | g \cdot x = x\} \) is trivial for all \( x \in M \).
  \item \textbf{Propery Discontinuously:}
    \begin{itemize}
      \item Any two points in different orbits have neighborhoods
        that do not intersect under the group action.
      \item for any point $ x $ the isotropy group $ G_x=\{g\in
        G|g\cdot x=x\} $ is finite.
      \item For any point $ x\in M $, there exists an open
        neighborhood $ U $ of $ x $ such that for all $ g\in G $ not
        in the isotropy group $ G_x $, the images $ g\cdot U $ do not
        intersect with $ U $, i.e., $ g\cdot U\cap U=\varnothing $.
        for all $ g \in G_x $, we have $ g\cdot U=U $.
    \end{itemize}
\end{itemize}
\defi{
  Orbifold

  If a group \( G \) acts properly discontinuously but not freely on
  a manifold \( M \), the quotient space \( M/G \) is called an
  \textbf{orbifold}.
}
An orbifold is a generalization of manifold with some mild singularities.

We can find a point on $ M $ to label the orbits in $ M/G $, the
complete set of such points is called a \textbf{fundamental region}
of the group action.

\subsubsection{Topology of Quotient Space}

\thm{
  Topology of Quotient Space of Simply Connected Manifold

  If $ M $ is a simply connected manifold and $ G $ is a finite group
  then we can prove a relation between the fundamental group of the
  quotient space and the group action:
  \begin{align}
    \pi_1(M/G)\cong G.
  \end{align}
}

In this case the projection $ p:M\to M/G $ is the universal covering
map of the quotient space and $ M $ is the universal covering space of $ M/G $.

\thm{
  Connected Manifold as Quotient Space

  Any connected manifold \( M \) can be represented as the quotient
  space of its universal covering space \( \tilde{M} \) by the action
  of its fundamental group \( \pi_1(M) \):
  \begin{align}
    M \cong \tilde{M}/\pi_1(M).
  \end{align}
  Here we define the action of \( \pi_1(M) \) on \( \tilde{M} \) as:
  Mapping the loop in $ M $ mapped to open path in $ \tilde{M} $ back
  by identifying its endpoints.
}

Mathematically, this action of \( \pi_1(M) \) on \( \tilde{M} \) is
defined as a \textbf{Deck Transformation}:
\defi{
  Deck Transformation

  For any $ [\gamma] \in \pi_1(M) $ we can define a map $
  D_{[\gamma]}: \tilde{M}\to\tilde{M} $ , that satisfy $ p\circ
  D_{[\gamma]}=p $. It gives a representation of \( \pi_1(M) \) as a
  group of homeomorphisms of \( \tilde{M} \).
}

\subsection{Geometization in 2D }

Before we only consider topological manifolds, now we can further
consider differential manifolds with riemannian metrics.

\subsubsection{Standard Geometries in 2D}

In 2 dimensions, there are three standard geometries with constant curvature:
\begin{itemize}
  \item \textbf{Riemann Sphere} $ \hat{\mathbb{C}} = \mathbb{C}\cup
    \infty $ with constant positive curvature metric:
    \begin{align}
      ds^2=\frac{4dzd\bar{z}}{(1+|z|^2)^2}.
    \end{align}
  \item \textbf{Complex Plane} $ \mathbb{C} $ with flat metric:
    \begin{align}
      ds^2=dzd\bar{z}.
    \end{align}
  \item \textbf{Hyperbolic Plane} $ \mathbb{H}^2 =
    \{z\in\mathbb{C}|\text{Im}(z)>0\} $ with constant negative curvature metric:
    \begin{align}
      ds^2=\frac{dzd\bar{z}}{(\text{Im}(z))^2}.
    \end{align}
\end{itemize}

\subsubsection{Uniformization Theorem}

We give out one of the most important theorems in 2D geometry:
\thm{\label{thm:Uniformization Theorem}
  Uniformization Theorem(Pure Topology Version)

  Every Surface is homeomorphic (diffeomporphic) to $ D/\Gamma $,
  where $ D \in \{\hat{\mathbb{C}}, \mathbb{C}, \mathbb{H}^2 \} $ and
  $ \Gamma\subset \text{Iso}(D) $ is a properly discontinuously
  isometry on $ D $.
}
This is a topological statement, we only say that a topological
surface can be represented as a quotient space of one of the three
standard geometries. However, with suitable operations we can give
the quotient space a constant curvature metric induced from the
standard geometry. Thus, we know that:
\begin{itemize}
  \item Any closed surface can be given a constant curvature metric.
  \item Classification of such metric = Classification of isometry
    group actions on the standard geometries.
\end{itemize}
